\chapter{Linear Operator Galois Correspondence}

Fix a Picard--Vessiot ring \(R/k\), its fundamental matrix \(Y\),
and its total extension \(Q=Q(R)\), under the hypotheses of the
previous chapter. In this chapter an automorphism of an operator
ring always commutes with the operator. The difference case requires
intermediate total rings because even the simplest nontrivial
example in the preceding chapter can have orthogonal idempotents.

\section{The algebra of the Galois group}

\begin{lemma}[Reduced operator tensors in characteristic zero]
\label{res:PV11}
If \(A\) is a reduced finite-type algebra over a characteristic-zero
field \(F\), and \(L/F\) is any field extension, then
\(L\otimes_F A\) is reduced. The same conclusion holds when the base
is a finite product of such fields, \(A\) is again reduced and
of finite type, and \(L\) is a finite product of field extensions
of the base components.
\end{lemma}

\begin{proof}
The embeddings of the domains \(A/P\), with \(P\) running over the
finitely many minimal primes, into their fraction fields embed \(A\)
in a finite product of finitely generated fields over \(F\).
Tensoring over a field preserves this injection. For one such
fraction field \(M\), choose a transcendence basis \(t\) of
\(M/F\). Then \(M/F(t)\) is finite separable, so it is generated
by a root of a monic separable polynomial \(p\) by the primitive
element theorem~\cite[Section 14.31, PDF p.~113]{mg-d07}. The algebra
\(L\otimes_F F(t)\) embeds in \(L(t)\), as it is the localization
of \(L[t]\) at the nonzero polynomials of \(F[t]\).
The algebra \(L\otimes_F M\) is obtained from that localization
by adjoining a root of \(p\). The basis of powers of the root
exhibits an injection into \(L(t)[T]/(p)\), which is reduced
as \(p\) remains separable. This proves the field case. Decomposing
the tensor product componentwise proves the product case.
\end{proof}

\begin{theorem}[The torsor of a fundamental matrix]
\label{res:PV12}
Write \(Y_1=Y\otimes1\), \(Y_2=1\otimes Y\) in \(R\otimes_k R\),
and let \(Z=Y_1^{-1}Y_2\). Then \(Z\) is constant and the
constant algebra \(D=C[Z_{ij},\det(Z)^{-1}]\) is reduced. There
is a natural isomorphism
\[
 R\otimes_C D \ \xrightarrow{\ \sim\ }\ R\otimes_k R
\]
by using \(R\) in the first factor. The evaluations \(D\to C\)
identify the matrices of \(Z\) with the algebraic subgroup
\[
 G=\{g\in GL_n(C):Y\mapsto Yg
                   \text{ gives an operator automorphism of }R/k\}.
\]
In particular \(D=C[G]\). Every such automorphism extends uniquely
to \(Q\), and every operator automorphism of \(Q/k\) restricts to
\(R\). Replacing \(Y\) by another fundamental matrix conjugates
the matrix realization of \(G\).
\end{theorem}

\begin{proof}
Constancy of \(Z\) follows from Lemma~\ref{res:PV02}. Since
\(Y_2=Y_1Z\), the right tensor product is generated by \(D\)
over its first factor, and Lemma~\ref{res:PV08} proves the
isomorphism. The embedding of \(R\) into the product of its
fraction fields and Lemma~\ref{res:PV11} show that
\(R\otimes_k R\) is reduced, hence so is \(D\).

An evaluation \(D\to C\) sends \(Z\) to an invertible matrix \(g\);
it maps the second copy of \(R\) to the first by \(Y\mapsto Yg\).
This is an operator homomorphism, injective by simplicity and
surjective since \(Y=(Yg)g^{-1}\). Conversely any automorphism
determines an evaluation by sending \(a\otimes b\) to \(a\,g(b)\),
where \(g(b)\) denotes its action. Its matrix is constant by
Lemma~\ref{res:PV02}. Thus \(G\) is the zero set in \(GL_n(C)\)
of the relations of \(Z\), it is a subgroup by composition,
and its coordinate algebra is exactly \(D\) by
Lemma~\ref{res:PV04}. The result is the torsor formula: in
\(R\otimes_k R\) the second matrix is obtained from the first by
right multiplication by \(Z\).

An automorphism of \(R\) preserves non-zero-divisors and
extends uniquely. For an automorphism of \(Q/k\), the ratio of its
image of \(Y\) and \(Y\) is constant. Since \(Q^{\mathrm{const}}=C\)
by Theorem~\ref{res:PV07}, the image of \(Y\) is of the form \(Yg\)
and generates \(R\) over \(k\). This proves restriction. For
another fundamental matrix \(Yb\), the new coordinate matrix of
an automorphism is \(b^{-1}gb\).
\end{proof}

\section{Intermediate rings and invariant elements}

\begin{lemma}[Structure of an intermediate total ring]
\label{res:PV13}
Let \(F\subseteq Q\) be an operator subring containing \(k\)
such that every non-zero-divisor of \(F\) is a unit of \(F\).
In the differential case \(F\) is a field. In the difference
case \(F\) is a finite product of fields with the operator permuting
its factors transitively. In either case \(F^{\mathrm{const}}=C\).
If \(R_F=F[Y_{ij},\det(Y)^{-1}]\) inside \(Q\), then \(R_F\)
is noetherian and its total ring of fractions is \(Q\).
\end{lemma}

\begin{proof}
If \(Q\) is a field, the subring \(F\) is a domain and the
unit condition makes it a field. For the difference case write
\(Q=\prod_i Q_i\) by Lemma~\ref{res:PV03}. The distinct kernels
of the projection maps \(F\to Q_i\) are permuted transitively.
No two distinct such kernels are comparable: any strict inclusion
would, by a power of the operator carrying one kernel to the other,
lead to a cycle of strict inclusions in a finite set.
For each distinct kernel \(P\) choose an element \(t_P\) that
belongs to every other kernel and does not belong to \(P\).
It exists as a product of elements supplied by incomparability.
Then \(s=\sum_P t_P\) has no zero coordinate in \(Q\), hence
is a non-zero-divisor of \(F\) and a unit there. The element
\(t_P/s\) is the characteristic idempotent of the set of
coordinates with kernel \(P\). This decomposes \(F\) into domains.
If an element in one domain component is nonzero, add the
identity of all other components to obtain a non-zero-divisor
of \(F\), then use its inverse to invert the original element
in that component. The components are therefore fields.
Constancy follows from \(k\subseteq F\subseteq Q\).

The algebra \(R_F\) is finite type over a finite product of
fields, hence noetherian. It contains the factor idempotents
of \(R\) supplied by Lemma~\ref{res:PV03}. Thus any zero
coordinate of an element of \(R_F\) makes it a zero divisor
unless \(Q\) has just one factor, when the same criterion
holds in the domain. Its non-zero-divisors are precisely the
elements invertible in \(Q\). Every fraction from \(R\) can
therefore be formed from \(R_F\), and every fraction from \(R_F\)
can be formed in \(Q\). The two total rings agree.
\end{proof}

\begin{theorem}[Fixed ring over an intermediate total ring]
\label{res:PV14}
For \(F\) as in Lemma~\ref{res:PV13}, let \(G_F\) consist of the
elements of \(G\) fixing \(F\) pointwise. Then \(G_F\) is an
algebraic subgroup of \(G\) and \(Q^{G_F}=F\).
\end{theorem}

\begin{proof}
The product \(Q\otimes_F R_F\), with the operator in both
factors, contains its first factor faithfully. In fact, on
each component of \(F\) both algebras contain the component
field, and tensoring over a field preserves injections.
It is generated over \(Q\) by the constant matrix of ratios
\(Z_F=(Y\otimes1)^{-1}(1\otimes Y)\) and its inverse determinant.
The ring \(Q\) is simple as an operator ring, since it is a
field or a transitive product of fields. Lemma~\ref{res:PV08}
therefore gives
\[
 Q\otimes_F R_F \ \cong\
 Q\otimes_C C[(Z_F)_{ij},\det(Z_F)^{-1}].
\]
The algebra on the left is reduced by
Lemma~\ref{res:PV11}, since \(R_F\subseteq Q\) is reduced and
Lemma~\ref{res:PV13} identifies the base components. Thus
the constant algebra on the right is reduced.

Its \(C\)-valued evaluations induce operator maps \(R_F\to Q\)
fixing \(F\) and sending \(Y\) to \(Yg\), \(g\in GL_n(C)\).
Such a map has image exactly \(R_F\) and is a surjective
endomorphism of a noetherian ring. It is injective: the ascending
sequence of kernels of its powers stabilizes, and surjectivity
then forces the first kernel to be zero. It extends to an
automorphism of \(Q/F\) by Lemma~\ref{res:PV13}; by
Theorem~\ref{res:PV12} it belongs to \(G_F\). Conversely every
element of \(G_F\) gives such an evaluation. This realizes \(G_F\)
as the algebraic subgroup of \(G\) cut out by the relations of
\(Z_F\).

Suppose that \(u=r/s\in Q\), with \(r,s\in R\) and \(s\)
a non-zero-divisor, is fixed by \(G_F\). The element
\(1\otimes r-u\otimes s\) of \(Q\otimes_F R_F\) vanishes under
every evaluation of the constant algebra. Express this element
as a finite sum with \(C\)-linearly independent coefficients in
\(Q\). Lemma~\ref{res:PV04} implies that each of its constant
algebra coefficients is zero. Its image in \(Q\otimes_F Q\)
is therefore zero. Since \(1\otimes s\) is invertible there,
\(1\otimes u=u\otimes1\).
To see that \(u\in F\), decompose over the field factors of \(F\)
and choose in each component an \(F\)-linear map from the
corresponding component of \(Q\) to that of \(F\), sending \(1\)
to \(1\). Applying this map in the first factor to the last
tensor equality gives \(u\) in the corresponding copy of \(F\).
This proves the fixed-ring claim.
\end{proof}

\section{Separation of subgroups}

\begin{lemma}[A representation detecting a closed subgroup]
\label{res:PV15}
For a closed subgroup \(H\subseteq G\), there exist a finite
dimensional regular representation \(\rho\) of \(G\) over \(C\)
and a line \(\ell\) in its space whose stabilizer is exactly \(H\).
If \(H\) is normal, there is a finite dimensional regular
linear representation of \(G\) with kernel exactly \(H\).
\end{lemma}

\begin{proof}
The ideal of functions vanishing on \(H\) in \(C[G]\) is
finitely generated. Take a common power of the determinant in
the denominators of representatives of the generators and bound
the degrees of all their numerators. The functions with those
bounds form a finite dimensional vector space \(V\) stable
under right translation, \(\rho_0(g)f(x)=f(xg)\).
The representation is regular: translating affects each bounded
monomial by polynomial coefficients in entries of \(g\) and
\(\det(g)^{-1}\). Taking a basis of \(V\) expresses the translation
in regular matrix coefficients. Let \(W\subseteq V\) consist of
the functions vanishing on \(H\). Every \(h\in H\) preserves
\(W\). Conversely, if \(g\) preserves \(W\), each of the chosen
ideal generators \(f\in W\) satisfies \(f(g)=(\rho_0(g)f)(1)=0\),
so \(g\in H\). If \(d=\dim W\), take
\(\rho=\bigwedge^d\rho_0\) and \(\ell=\bigwedge^d W\).
For a nonzero \(w\in\ell\) the equality \(v\wedge w=0\) in
\(\bigwedge^{d+1}V\) characterizes \(v\in W\), including when
\(W=0\) or \(W=V\). Thus the stabilizer of \(\ell\) is \(H\).

Assume now \(H\) is normal. For each irreducible component
\(X\) of \(G\), let \(U_X\) be the span of the vectors
\(\rho(g)w\) with \(g\in X\). For \(h\in H\), normality
makes each such vector an eigenvector of \(\rho(h)\).
The locus where the eigenvalue is a given scalar is closed
in \(X\), and there are only finitely many eigenvalues.
Irreducibility forces one eigenvalue throughout \(X\).
Hence \(h\) acts as a scalar on \(U_X\).
Left translation by \(g\) permutes the components, so
\(\rho(g)\) permutes the finite set of distinct spaces \(U_X\).
On the direct sum of the spaces \(\operatorname{End}_C(U_X)\),
one summand for every distinct \(U_X\), act by transport and
conjugation using \(\rho(g)\). This is regular. For each possible
permutation of the summands, the set realizing it is
closed, as transport of specified subspaces is given by linear
conditions on the matrix. Its complement is a finite union of such
closed sets, so the locus is open as well. Functions regular on
these finitely many closed-and-open pieces combine into regular
functions on \(G\): the vanishing ideals of distinct pieces are
comaximal by Lemma~\ref{res:PV04}, and their intersection is zero.
For two comaximal ideals \(I,J\), write \(1=i+j\) with \(i\in I,j\in J\).
The element \(aj+bi\) then combines \(a\bmod I\) and \(b\bmod J\).
Iterating gives the finite decomposition.
On each piece, the formulas come from the regular matrices of
\(\rho(g)\) and \(\rho(g^{-1})\).
This action kills \(H\). If it kills \(g\), then \(g\)
preserves every summand and acts by a scalar on each \(U_X\):
commuting with coordinate projections makes a matrix diagonal,
and commuting with all matrix units makes the diagonal entries equal.
One of the \(U_X\) contains \(w\), so \(g\) preserves \(\ell\)
and belongs to \(H\). This is the asserted representation.
\end{proof}

\begin{lemma}[Separating a coset by a rational invariant]
\label{res:PV16}
Suppose \(H\subseteq G\) is closed and \(g\in G\setminus H\).
Then \(g\) acts nontrivially on \(Q^H\).
\end{lemma}

\begin{proof}
First consider the coordinate algebra \(C[G]\) and its total
ring of fractions. Take \(\rho,\ell=Cw\) as in
Lemma~\ref{res:PV15}. Choose a linear functional \(a\)
nonzero on \(w\), on \(\rho(g)w\), and on at least one vector
\(\rho(x)w\) from each irreducible component of \(G\).
This is possible over the infinite field \(C\) by avoiding
finitely many proper hyperplanes. Thus \(a(\rho(x)w)\)
is a non-zero-divisor in \(C[G]\). Since the lines of
\(w\) and \(\rho(g)w\) are distinct, choose a linear functional
\(b\) that takes different values on their vectors normalized
by \(a=1\). The rational function
\[
 f(x)=\frac{b(\rho(x)w)}{a(\rho(x)w)}
\]
is invariant under right translations by \(H\), because those
translations multiply numerator and denominator by the same
nonzero scalar. It has distinct values at \(1\) and \(g\).
Consequently right translation by \(g\) does not fix \(f\).

We transport this separation from \(G\) to \(Q\). Choose a
maximal ideal in \(R\) and let \(E\) be its residue field.
It is a finite extension of \(k\) by Lemma~\ref{res:PV04}.
Base-changing the first factor of Theorem~\ref{res:PV12}
along \(R\to E\) gives
\[
 E\otimes_k R\cong E\otimes_C C[G].
\]
Here an element \(g\in G\) acts on the first algebra by its
action on \(R\) and on the second by right translation,
because \(Y_1^{-1}Y_2\) changes to \(Y_1^{-1}Y_2g\).
The total ring of fractions on the left is \(E\otimes_k Q\).
Indeed this tensor is a product of fields, since
\(E/k\) is finite separable and each factor of \(Q\) is a field.
It is a localization of \(E\otimes_k R\), at non-zero-divisors
from \(R\). Any non-zero-divisor of \(E\otimes_k R\) remains
a non-zero-divisor after localization and hence is a unit
in the resulting product of fields. On the right, the
denominator of \(f\) stays a non-zero-divisor after extension
from \(C\) to \(E\), by exactness of tensoring over a field.
Thus \(f\) gives an \(H\)-invariant element of \(E\otimes_k Q\)
moved by \(g\); the distinct evaluations at \(1,g\) persist
over \(E\). Since invariants of \(E\otimes_k Q\) are
\(E\otimes_k Q^H\), as one sees by expanding in a \(k\)-basis
of \(E\), the action of \(g\) on \(Q^H\) is nontrivial.
\end{proof}

\begin{theorem}[Differential and difference Galois correspondence]
\label{res:PV17}
The maps
\[
 H\longmapsto Q^H,\qquad
 F\longmapsto G_F
\]
are inverse inclusion-reversing bijections between closed
subgroups of \(G\) and intermediate operator subrings \(F\)
in which every non-zero-divisor is a unit. In the differential
case the latter are precisely the intermediate differential
fields. In particular \(Q^G=k\).
\end{theorem}

\begin{proof}
For any subgroup \(H\), its fixed subring is stable under the
operator and contains \(k\). If an element \(u\in Q^H\) vanishes
in some component of \(Q\), the characteristic idempotent
of its zero components is nonzero, belongs to \(Q^H\), and
annihilates \(u\). Indeed every element of \(H\) permutes
the field factors and preserves this zero-component set.
Thus a non-zero-divisor of \(Q^H\) has no zero component;
its inverse in \(Q\) is fixed by \(H\). This proves that
\(Q^H\) is an intermediate ring of the indicated kind.
For such a ring \(F\), Theorem~\ref{res:PV14} gives \(Q^{G_F}=F\).
For a closed subgroup \(H\), each element of \(G\setminus H\)
is excluded from \(G_{Q^H}\) by Lemma~\ref{res:PV16}, and
\(H\subseteq G_{Q^H}\) by definition. The inclusion reversal
is immediate from the two definitions. Taking \(F=k\) gives
the full fixed-ring equality.
\end{proof}

\section{Normal subgroups and Picard--Vessiot quotients}

\begin{lemma}[Finite semilinear descent of a matrix]
\label{res:PV18}
Over a field \(k\) of arbitrary characteristic, suppose \(E/k\) is
finite Galois, with group \(\Gamma\), and \(a_\gamma\in GL_m(E)\) obeys
\(a_{\gamma\eta}=a_\gamma\gamma(a_\eta)\).
Then for some \(B\in GL_m(E)\) one has
\(\gamma(B)=B a_\gamma\) for all \(\gamma\in\Gamma\).
\end{lemma}

\begin{proof}
On row vectors in \(E^m\) put \(T_\gamma(v)=\gamma(v)a_\gamma^{-1}\).
The cocycle identity makes this an action by semilinear maps.
The sums \(P(v)=\sum_\gamma T_\gamma(v)\) are invariant.
They span \(E^m\) over \(E\): if a column \(c\) annihilates all
of them, substitute \(v=x e_i\) for each coordinate row \(e_i\).
Then \(\sum_\gamma \gamma(x)e_i a_\gamma^{-1}c=0\) for all \(x\).
Distinct field automorphisms are linearly independent as functions
on \(E\) by~\cite[Lemma 16.6, PDF p.~125]{mg-d07}.
Thus the coefficient of the identity automorphism gives
\(e_i c=0\) for all \(i\).
Taking \(m\) independent invariant rows gives \(B\) as required.
\end{proof}

\begin{theorem}[The normal subgroup criterion and quotient]
\label{res:PV19}
An intermediate total ring \(F=Q^H\) of Theorem~\ref{res:PV17}
is stable under the action of \(G\) if and only if \(H\)
is normal. In that case \(F\) is the total extension of a
Picard--Vessiot ring for a system over \(k\), and restriction
is a regular morphism in the matrix realizations and gives an
isomorphism of groups of \(C\)-points
\[
 G/H \ \cong\ \operatorname{Aut}_{\mathrm{op}}(F/k).
\]
Conversely, any intermediate total ring that is a
Picard--Vessiot total extension over \(k\) is stable under \(G\).
\end{theorem}

\begin{proof}
The correspondence of Theorem~\ref{res:PV17} and
\(g(Q^H)=Q^{gHg^{-1}}\) give the stability criterion.
If \(H\) is normal, take a regular representation
\(\rho:G\to GL_m(C)\) with kernel \(H\) from
Lemma~\ref{res:PV15}. There is a finite Galois extension \(E/k\)
over which the torsor has an \(E\)-point: take a finite residue
extension as in Lemma~\ref{res:PV16} and then a splitting
field of its separable defining polynomial. If \(U\in GL_n(E)\)
denotes the image of \(Y\) at this point, the torsor coordinates
on \(E\otimes_k R\) are \(Z=U^{-1}Y\). Writing
\(\gamma(U)=U c_\gamma\), the torsor formula puts
\(c_\gamma\in G(E)\) and gives
\(c_{\gamma\eta}=c_\gamma\gamma(c_\eta)\). Here \(G(E)\)
means the \(E\)-valued solutions of the relations for \(G\).
The multiplication law and the regular representation extend
to points valued in any \(C\)-algebra. Indeed their polynomial identities on
\(C\)-points are identities in the coordinate algebras,
by Lemma~\ref{res:PV04}, also for products. For the latter
assertion, successive evaluation in each factor of a
tensor product of coordinate algebras detects any nonzero
element by independence of coefficients.

By Lemma~\ref{res:PV18} take \(B\) with
\(\gamma(B)=B\rho(c_\gamma)\). Then
\[
 Y_H=B\rho(Z)\in GL_m(E\otimes_k R)
\]
is invariant under \(\Gamma\), since
\(\gamma(Z)=c_\gamma^{-1} Z\).
Expanding tensor elements in a \(k\)-basis of \(R\) proves
that \((E\otimes_k R)^\Gamma=R\), so \(Y_H\) and its inverse
are in \(R\). The action of an element \(g\in G\) on this
matrix is \(Y_H\mapsto Y_H\rho(g)\). Its entries are therefore
fixed by \(H\). The coefficient matrix
\(\partial(Y_H)Y_H^{-1}\) or \(\sigma(Y_H)Y_H^{-1}\) is
fixed by all of \(G\), as the right factors \(\rho(g)\)
are constant. By Theorem~\ref{res:PV17} it lies in \(M_m(k)\);
in the difference case it is invertible there, by applying
the same argument to its inverse. Therefore \(Y_H\) is a
fundamental matrix for a system over \(k\).

Its ring \(R_H=k[(Y_H)_{ij},\det(Y_H)^{-1}]\) is
Picard--Vessiot by Theorem~\ref{res:PV09}, applied to \(Q\).
Each non-zero-divisor of \(R_H\) is invertible in \(Q\).
For the difference case, use the distinct projection kernels
from \(R_H\) to the field factors of \(Q\). They are
transitively permuted and hence incomparable, as in the proof
of Lemma~\ref{res:PV13}. If an element vanishes in one
projection, it annihilates the nonzero product belonging to
all other distinct kernels and outside that kernel. Thus
every non-zero-divisor has nonzero image in each factor.
It follows that \(Q(R_H)\) is an intermediate total ring in
\(Q^H\). An automorphism fixes \(Q(R_H)\) precisely if it
fixes \(Y_H\), or equivalently if it belongs to \(H\).
Theorem~\ref{res:PV17} now yields \(Q(R_H)=F\).

In the matrix realizations restriction is \(g\mapsto\rho(g)\),
so it is regular and its kernel is \(H\).
To prove surjectivity, let \(\tau\) be an operator
automorphism of \(F/k\). Form \(Q\otimes_F R_F\) with the
map \(F\to Q\) on the first side twisted by \(\tau\).
It is nonzero and contains \(Q\) by the field-component
argument of Theorem~\ref{res:PV14}. It is generated over
\(Q\) by the same constant ratios of fundamental matrices.
Lemma~\ref{res:PV08} and a \(C\)-point of their constant
algebra give a map \(R_F\to Q\) extending \(\tau\) and sending
\(Y\) to \(Yg\) for some \(g\in GL_n(C)\). Its image is
\(R_F\); it is an isomorphism by the noetherian kernel argument
in Theorem~\ref{res:PV14}. It extends to \(Q\), yielding
the desired lift in \(G\). Finally, if an intermediate ring is
given as a Picard--Vessiot total extension, a fundamental
matrix for it is carried by every \(g\in G\) to another
fundamental matrix for the same system. Its constant ratio
lies in \(GL_m(C)\). The ring of entries and its total ring
are preserved, proving the converse.
\end{proof}

\begin{corollary}[The algebraic part of a differential extension]
\label{res:PV20}
In the differential case, let \(G^0\) be the connected component
of the identity in the Zariski topology of \(G\). Then \(Q^{G^0}\)
is the relative algebraic closure of \(k\) in \(Q\),
and it is a finite Galois extension of \(k\) with group \(G/G^0\).
\end{corollary}

\begin{proof}
The component \(G^0\) is a closed normal subgroup.
Indeed inversion and conjugation preserve the identity component,
and for \(g\in G^0\) the connected set \(gG^0\) meets \(G^0\) at
\(g\), hence lies in \(G^0\). The components are its translates, and
there are finitely many since each is a union of irreducible
components of a noetherian space.
For \(u\in Q\) algebraic over \(k\), its orbit is finite because
all its translates are roots of its minimal polynomial.
Its stabilizer is the closed subgroup \(G_{k(u)}\) by
Theorem~\ref{res:PV14}. A closed subgroup of finite index is also
open because its complement is a finite union of its other cosets.
Thus the stabilizer contains \(G^0\), proving \(u\in Q^{G^0}\).

Conversely, every \(u\in Q^{G^0}\) has finite \(G\)-orbit of
cardinality at most \(d=|G/G^0|\). The polynomial whose roots are its
distinct translates has coefficients in \(Q^G=k\) by
Theorem~\ref{res:PV17}. It splits in \(Q^{G^0}\) and is separable.
Thus each \(u\) is algebraic of degree at most \(d\), and its minimal
polynomial splits in \(Q^{G^0}\). Any finite subextension is
separable and has a primitive element by
\cite[Section 14.31, PDF p.~113]{mg-d07}, so also has degree at most
\(d\). A finite subextension of largest degree must contain
every \(u\), since adjoining any further element would increase the
degree. Hence \(Q^{G^0}/k\) is finite, normal, and separable. The
automorphisms in Theorem~\ref{res:PV19} are all its field
automorphisms, since a derivation extends uniquely to a finite
separable field extension. Uniqueness follows by differentiating
the minimal polynomial of each element, whose derivative at that
element is nonzero.
\end{proof}

\section{Two rank-one calculations}

\begin{corollary}[The rank-one multiplicative group]
\label{res:PV21}
Consider \(\partial y=a y\), \(a\in k\), or
\(\sigma(y)=a y\), \(a\in k^\times\). The Galois group is
either \(C^\times\) or the group \(\mu_m(C)\) of \(m\)-th roots
of unity. The finite case occurs exactly when there is some
positive \(m\) and \(b\in k^\times\) with
\[
 \partial b=ma b \quad\text{or}\quad \sigma(b)=a^m b,
\]
respectively. The least such \(m\) is the order of the group.
\end{corollary}

\begin{proof}
Every element \(g\in G\) acts on a fundamental unit \(y\) by
\(y\mapsto gy\). By Theorem~\ref{res:PV12}, \(G\) is a
closed subgroup of \(C^\times\), so it is finite unless it is
all of \(C^\times\), since a nonzero Laurent polynomial has
finitely many roots. A finite subgroup of order \(m\) lies
in \(\mu_m(C)\) by the group order theorem, hence equals it.
The element \(y^m\) is fixed exactly when \(G\subseteq\mu_m(C)\),
so by Theorem~\ref{res:PV17} this condition is equivalent to
\(y^m\in k^\times\). Then \(b=y^m\) satisfies the required
equation. Conversely, \(y^m/b\) is constant for any such
\(b\) and belongs to \(C^\times\), hence \(y^m\in k^\times\).
This proves both the criterion and its minimality assertion.
\end{proof}

\begin{example}[A disconnected difference Galois group]
\label{res:PV22}
The extension of Example~\ref{res:PV10} has Galois group
\(\{1,-1\}\), acting by \(u\mapsto\pm u\), and its full fixed
ring is \(k\). For the differential field \(C(z)\) with
\(\partial z=1\), the equation \(\partial y=y\) has Galois
group \(C^\times\).
\end{example}

\begin{proof}
In Example~\ref{res:PV10}, the relation \(u^2=1\) forces the
multiplier to have square \(1\); both choices give
automorphisms. The nonidentity automorphism interchanges the
two factors of \(k\times k\) and fixes its diagonal.
In the differential example the constants are \(C\), since
a rational function of derivative zero in characteristic zero
cannot have a finite pole and then must be a constant polynomial.
For any nonzero relatively prime polynomials \(p,q\), the
logarithmic derivative of \(p/q\) is
\((p'q-pq')/(pq)\), whose numerator has smaller degree than
its denominator whenever it is nonzero. It cannot equal
a positive integer \(m\) in characteristic zero, so
Corollary~\ref{res:PV21} gives the claim.
\end{proof}

\paragraph{Sources.}
For differential Galois groups and the correspondence see
van der Put--Singer~\cite[Chapter 1, Section 1.4]{vps-differential}
and Singer~\cite[Sections 1.3 and 1.5]{singer}. For the total
ring and difference correspondence see van der Put--Singer
\cite[Chapter 1, Sections 1.2--1.3]{vps-difference} and
Roques~\cite[Sections 5.2--8]{roques}.
